%% Question bank chapter file - generated by bank_tools.py build (do not edit while the bank is still being built:
%% it is regenerated from the per-test data in Question Bank\data). Every question is one bankq block:
%%   \begin{bankq}[id=..., source={...}, paper=.., qno=.., type=..., subtopic={...}, answer={...}, answer_src=key|solved,
%%                 working={...}, difficulty=1-10, diff_parts={K S I C F}, diff_note={...},
%%                 group=<passage id or empty>, also_in={duplicate copies in other tests}]
%% difficulty = min(10, K+S+I+C+F-1): K concepts 1-3, S steps 1-4, I insight 0-2, C computation 0-2, F format 0-1
%%   ...question LaTeX, options, figures (\bfig / \bfigi -> img/)...
%%   \end{bankq}
%% Passages shared by several questions are bankpassage blocks; their questions carry group=<passage id>.
\documentclass[10pt]{article}
\usepackage[a4paper,margin=1.8cm]{geometry}
\usepackage[T1]{fontenc}
\usepackage{amsmath,amssymb,mathtools}
\usepackage{graphicx}
\usepackage[export]{adjustbox}
\usepackage{xcolor,needspace,keyval}
\usepackage{enumitem}
\usepackage{booktabs,array,tabularx}
\usepackage[hidelinks]{hyperref}
\setlength{\parindent}{0pt}
\setlength{\parskip}{3pt}
\makeatletter
\define@key{bq}{id}{\def\bq@id{#1}}\define@key{bq}{source}{\def\bq@src{#1}}\define@key{bq}{paper}{\def\bq@paper{#1}}
\define@key{bq}{qno}{\def\bq@qno{#1}}\define@key{bq}{type}{\def\bq@type{#1}}\define@key{bq}{subtopic}{\def\bq@sub{#1}}
\define@key{bq}{answer}{\def\bq@ans{#1}}\define@key{bq}{answer_src}{\def\bq@asrc{#1}}\define@key{bq}{working}{\def\bq@work{#1}}
\define@key{bq}{group}{\def\bq@group{#1}}\define@key{bq}{also_in}{\def\bq@also{#1}}
\define@key{bq}{difficulty}{\def\bq@diff{#1}}\define@key{bq}{diff_parts}{\def\bq@dparts{#1}}\define@key{bq}{diff_note}{\def\bq@dnote{#1}}
\newenvironment{bankq}[1][]{\def\bq@work{}\def\bq@also{}\def\bq@group{}\def\bq@diff{?}\def\bq@dparts{}\def\bq@dnote{}\setkeys{bq}{#1}%
  \par\needspace{4\baselineskip}\medskip\noindent{\small\textbf{\bq@id}\hfill\color{gray}[\bq@src, Paper \bq@paper, Q\bq@qno]
  \bq@type\ $\cdot$ \bq@sub\ $\cdot$ \textbf{D\bq@diff}}\par\nopagebreak\relax}
  {\par{\small\color{teal}Answer: \textbf{\bq@ans}\ifx\bq@work\@empty\else\ (solved: \bq@work)\fi
  \ifx\bq@also\@empty\else\par\color{gray}Also in: \bq@also\fi}\par}
\define@key{bp}{id}{\def\bp@id{#1}}\define@key{bp}{source}{\def\bp@src{#1}}\define@key{bp}{questions}{\def\bp@qs{#1}}
\newenvironment{bankpassage}[1][]{\setkeys{bp}{#1}\par\needspace{4\baselineskip}\bigskip\noindent
  {\small\textbf{Passage \bp@id}\hfill\color{gray}[\bp@src; questions \bp@qs]}\par\nopagebreak\relax\itshape}{\par}
\makeatother
\newcommand{\bfig}[2]{\par\smallskip{\centering\includegraphics[width=#1pt,max width=0.62\linewidth,max height=0.32\textheight]{img/#2}\par}\smallskip}
\newcommand{\bfigi}[2]{\includegraphics[width=#1pt,max width=\linewidth,max height=0.32\textheight,valign=c]{img/#2}}
\newcommand{\bmissing}[1]{\textit{[figure missing in the source conversion: \detokenize{#1}]}}
\newenvironment{qoptions}{\begin{itemize}[leftmargin=1.8em,itemsep=1pt,topsep=2pt,label={}]}{\end{itemize}}
\newcommand{\opt}[2]{\item[(#1)] #2}
\providecommand{\fwstart}{}\providecommand{\fwend}{}\providecommand{\qfullsplit}{}
\newcolumntype{L}[1]{>{\raggedright\arraybackslash}p{#1}}
\newcolumntype{C}[1]{>{\centering\arraybackslash}p{#1}}
\newcolumntype{R}[1]{>{\raggedleft\arraybackslash}p{#1}}
\providecommand{\sectionhdr}[1]{}\providecommand{\subsectionhdr}[1]{}


\begin{document}
{\Large\bfseries Solid State --- Question Bank}\par\smallskip
{\small Chemistry chapter C12. Questions from all converted test papers, grouped by source series. Answers: printed key where it matches the paper, otherwise solved (answer\_src=solved, with working). Exact repeats across papers appear once, with also\_in listing the other copies.}\par\medskip\hrule

\section*{Score-I papers}
\begin{bankq}[id=SI12-P2-CHEM-S2-Q04, source={Score-I - Test 12 (09 Mar 2022)}, paper=2, qno=4, type=numerical, subtopic={Unit cells}, answer={6.91}, answer_src=solved, working={pOH $=13-8=5$; $[\mathrm{OH^-}]^2=K_hC$ with $K_h=10^{-8}$ gives $C=0.01$ M, $M=129.6$; $a=500$ pm, $\rho=4M/(N_Aa^3)=6.91$ g/cm$^3$}, difficulty=7, diff_parts={K3 S3 I0 C1 F1}, diff_note={Hydrolysis pH formula gives C=0.01 M and M=129.6; rock salt Z=4, a=500 pm, density}]
Ionic solid Na$^+$A$^-$ crystallise in rock salt type structure. 2.592 gm of ionic solid salt NaA dissolved in water to make 2 litre solution. The pH of this solution is 8. If distance between cation and anion is 250 pm. Calculate density of ionic solid (in gm/cm$^3$).\\
{(Given: p$K_w=13$, p$K_a$(HA)$=5$, $N_A=6\times10^{23}$)}
\end{bankq}

\section*{TOAS papers}
\begin{bankq}[id=TOAS01-P1-CHEM-S1-Q01, source={TOAS - Test 01 (30 Jan 2022)}, paper=1, qno=1, type=single, subtopic={Defects}, answer={B}, answer_src=key, difficulty=1, diff_parts={K1 S1 I0 C0 F0}, diff_note={Match defect type with example (Frenkel/F-centre/FeO)}]
Match the types of defect given in Column-I with the statement given in Column-II.
\par\nobreak
{\small\setlength{\tabcolsep}{3pt}\begin{center}
\begin{tabular}{|c|L{1.9cm}|c|L{3.1cm}|}
\hline
\multicolumn{2}{|c|}{\textbf{Column-I}} & \multicolumn{2}{c|}{\textbf{Column-II}}\\
\hline
(i) & Impurity defect & (a) & NaCl with anionic sites called F-centres\\
\hline
(ii) & Metal excess defect & (b) & FeO doped with Fe$^{3+}$\\
\hline
(iii) & Metal deficiency defect & (c) & NaCl doped with Sr$^{2+}$ and some cationic sites vacant\\
\hline
\end{tabular}\end{center}}
\begin{qoptions}
\opt{A}{(i)-a, (ii)-b, (iii)-c}
\opt{B}{(i)-c, (ii)-a, (iii)-b}
\opt{C}{(i)-a, (ii)-c, (iii)-b}
\opt{D}{(i)-b, (ii)-c, (iii)-a}
\end{qoptions}
\end{bankq}
\begin{bankq}[id=TOAS01-P1-CHEM-S2-Q05, source={TOAS - Test 01 (30 Jan 2022)}, paper=1, qno=5, type=multi, subtopic={Unit cells}, answer={A,C,D}, answer_src=key, difficulty=8, diff_parts={K2 S4 I1 C1 F1}, diff_note={Counting atoms per cell (P1 Q3 R4 S3) and after removal of atoms on given lines, four options}]
A cubical crystal is such that the corners atoms are occupied by P atoms, face centers are occupied by Q atoms, alternate tetrahedral voids are occupied by R atoms and edge centers are occupied by S atoms. Identify the correct option(s) :
\begin{qoptions}
\opt{A}{If there are no defects in the crystal, then the formula should be $\mathrm{PQ_3S_4R_3}$}
\opt{B}{If all the atoms lying on one of the body diagonal are missing then formula should be $\mathrm{PQ_4,\ R_{8/3}S_4}$}
\opt{C}{If all the atoms lying on a line passing through edge centers and face-centre of a face are missing, then formula should be $\mathrm{P_2Q_5R_8S_5}$}
\opt{D}{If all the atoms lying on a line passing through body center and connecting opposite face centers are missing, then formula will be $\mathrm{PQ_2R_4S_3}$}
\end{qoptions}
\end{bankq}
\begin{bankq}[id=TOAS01-P1-CHEM-S2-Q06, source={TOAS - Test 01 (30 Jan 2022)}, paper=1, qno=6, type=multi, subtopic={Unit cells}, answer={A,B,C}, answer_src=key, difficulty=3, diff_parts={K1 S2 I0 C0 F1}, diff_note={Four recall statements on solids and crystal systems}]
Which of the following statements are correct -
\begin{qoptions}
\opt{A}{Amorphous solids are also called pseudo solids or supercooled liquids.}
\opt{B}{If a solid is very hard, and it is electrical insulator in solid as well as in molten state and melts at extremely high temperature, then it can be covalent or network solid}
\opt{C}{In Rhombohedral crystal system, only primitive unit cell is possible}
\opt{D}{In fluorite structure coordination number of cation is 4.}
\end{qoptions}
\end{bankq}
\begin{bankpassage}[id=TOAS01-P1-CHEM-S3-G1-2, source={TOAS - Test 01 (30 Jan 2022)}, questions={1-2}]
A hypothetical metal M (Atomic weight $=450$) crystallizes as a cubic crystal such that metal atoms touch each other along body diagonal of cube. Edge length of cubic unit cell is 1nm.\\
If this metal is melted then it's density in molten form is 1.2 gm/cc. ($N_A=6\times10^{23}$)
\end{bankpassage}
\begin{bankq}[id=TOAS01-P1-CHEM-S3-Q01, source={TOAS - Test 01 (30 Jan 2022)}, paper=1, qno=1, type=numerical, subtopic={Unit cells}, answer={1.50}, answer_src=key, difficulty=4, diff_parts={K2 S2 I0 C0 F1}, diff_note={Metal touching along body diagonal is bcc, rho=Z M/(N_A a^3)}, group=TOAS01-P1-CHEM-S3-G1-2]
Density of solid metal in gm/cc is -
\end{bankq}
\begin{bankq}[id=TOAS01-P1-CHEM-S3-Q02, source={TOAS - Test 01 (30 Jan 2022)}, paper=1, qno=2, type=numerical, subtopic={Unit cells}, answer={45.60}, answer_src=key, difficulty=7, diff_parts={K2 S3 I1 C1 F1}, diff_note={Atom volume from r=sqrt3 a/4, atoms per cc from molten density, volume fraction}, group=TOAS01-P1-CHEM-S3-G1-2]
Percentage of empty space in molten state of metal will be - \ [Use : $\pi\sqrt3=5.44$]
\end{bankq}
\begin{bankq}[id=TOAS01-P1-CHEM-S4-Q07, source={TOAS - Test 01 (30 Jan 2022)}, paper=1, qno=7, type=numerical, subtopic={Unit cells}, answer={2}, answer_src=key, difficulty=5, diff_parts={K1 S3 I0 C1 F1}, diff_note={Mass of fcc cell from density and edge, then atoms in 32 g}]
A crystalline solid of a pure substance has a face-centred cubic structure with a cell edge of 200 pm. If the density of the substance in the crystal is 8g $\mathrm{cm^{-3}}$, then the number of atoms present in 32g of the crystal is $N\times10^{24}$. The value of $N$ is
\end{bankq}
\begin{bankq}[id=TOAS01-P1-CHEM-S4-Q08, source={TOAS - Test 01 (30 Jan 2022)}, paper=1, qno=8, type=numerical, subtopic={Packing}, answer={8}, answer_src=key, difficulty=2, diff_parts={K1 S1 I0 C0 F1}, diff_note={Radius ratio 0.837 gives cubic hole, CN 8}]
If radius of $\mathrm{M^+}$ is 113 pm and radius of $\mathrm{X^-}$ is 135 pm then the expected value of coordination number of $\mathrm{X^-}$ in MX(s) is ?
\end{bankq}
\begin{bankq}[id=TOAS01-P2-CHEM-S1-Q01, source={TOAS - Test 01 (30 Jan 2022)}, paper=2, qno=1, type=single, subtopic={Packing}, answer={C}, answer_src=key, difficulty=1, diff_parts={K1 S1 I0 C0 F0}, diff_note={Radius ratio 0.52 lies in the octahedral range}]
If the value of ionic radius ratio $\left(\dfrac{r^+}{r^-}\right)$ is 0.52 in an ionic compound, the geometrical arrangement of ions in crystal is
\begin{qoptions}
\opt{A}{Tetrahedral}
\opt{B}{Planar}
\opt{C}{Octahedral}
\opt{D}{Pyramidal}
\end{qoptions}
\end{bankq}
\begin{bankq}[id=TOAS01-P2-CHEM-S1-Q02, source={TOAS - Test 01 (30 Jan 2022)}, paper=2, qno=2, type=single, subtopic={Packing}, answer={B}, answer_src=key, difficulty=2, diff_parts={K1 S1 I1 C0 F0}, diff_note={Square layers stacked directly (simple cubic): vacant 100(1-pi/6)}]
Identical 4 spheres are taken and are arranged in a layer of square packing touching each other as shown
\bfig{49.7}{TOAS01_p2_chem_s1_q02_1.png}
The percentage of vacant space is (in 3-D)
\begin{qoptions}
\opt{A}{$100\left(1-\dfrac{3\pi}{8}\right)$}
\opt{B}{$100\left(1-\dfrac\pi6\right)$}
\opt{C}{$100-\dfrac{3\pi}{8}$}
\opt{D}{$\dfrac\pi6$}
\end{qoptions}
\end{bankq}
\begin{bankq}[id=TOAS01-P2-CHEM-S1-Q03, source={TOAS - Test 01 (30 Jan 2022)}, paper=2, qno=3, type=single, subtopic={Unit cells}, answer={D}, answer_src=key, difficulty=2, diff_parts={K1 S2 I0 C0 F0}, diff_note={a=2(r+ + r-) and Na-Na distance a/sqrt2}]
If $r_{\mathrm{Na^+}}=95$ pm and $r_{\mathrm{Cl^-}}=181$ pm in NaCl (rock salt) structure. What is the shortest distance between $\mathrm{Na^+}$ ions ?
\begin{qoptions}
\opt{A}{778 pm}
\opt{B}{276 pm}
\opt{C}{195 pm}
\opt{D}{390 pm}
\end{qoptions}
\end{bankq}
\begin{bankq}[id=TOAS01-P2-CHEM-S1-Q04, source={TOAS - Test 01 (30 Jan 2022)}, paper=2, qno=4, type=single, subtopic={Packing}, answer={C}, answer_src=key, difficulty=5, diff_parts={K2 S3 I1 C0 F0}, diff_note={Locate tetrahedral and octahedral voids on the given plane of the ccp cell}]
In a hypothetical solid `C' atoms are found to form cubical close packed lattice. `A' atoms occupy all tetrahedral voids \& `B' atoms occupy all octahedral voids. `A' and `B' atoms
\bfig{54.6}{TOAS01_p2_chem_s1_q04_1.png}
are of appropriate size i.e. there is no distortion in ccp lattice of C atoms. Now if a plane as shown in the following figure is cut, then the cross section of this plane will look like.
\begin{qoptions}
\opt{A}{\bfigi{39.7}{TOAS01_p2_chem_s1_q04_A.png}}
\opt{B}{\bfigi{42.6}{TOAS01_p2_chem_s1_q04_B.png}}
\opt{C}{\bfigi{39.7}{TOAS01_p2_chem_s1_q04_C.png}}
\opt{D}{\bfigi{39.9}{TOAS01_p2_chem_s1_q04_D.png}}
\end{qoptions}
\end{bankq}
\begin{bankq}[id=TOAS01-P2-CHEM-S1-Q06, source={TOAS - Test 01 (30 Jan 2022)}, paper=2, qno=6, type=single, subtopic={Packing}, answer={B}, answer_src=key, difficulty=1, diff_parts={K1 S1 I0 C0 F0}, diff_note={One octahedral void per atom, 2/3 filled gives A3B2}]
In an ionic solid atoms `A' form hexagonal close packing crystal with atoms `B' in ($2/3^{\mathrm{rd}}$) octahedral voids. Empirical formula of solid is -
\begin{qoptions}
\opt{A}{AB}
\opt{B}{$\mathrm{A_3B_2}$}
\opt{C}{$\mathrm{A_2B}$}
\opt{D}{$\mathrm{A_2B_3}$}
\end{qoptions}
\end{bankq}
\begin{bankq}[id=TOAS01-P2-CHEM-S2-Q02, source={TOAS - Test 01 (30 Jan 2022)}, paper=2, qno=2, type=numerical, subtopic={Packing}, answer={4}, answer_src=key, difficulty=4, diff_parts={K2 S2 I0 C0 F1}, diff_note={Ratio 0.858 gives CsCl type, body diagonal sqrt3 a=2(r+ + r-)}, also_in={TOAS - Test 03 (27 Feb 2022) P2 Q3}]
In an ionic solid $r_{(+)}=1.6$ \AA\ and $r_{(-)}=1.864$ \AA. Use the radius ratio rule to determine the edge length of the cubic unit cell in \AA. [ $\sqrt3=1.732$ ]
\end{bankq}
\begin{bankq}[id=TOAS01-P2-CHEM-S2-Q04, source={TOAS - Test 01 (30 Jan 2022)}, paper=2, qno=4, type=numerical, subtopic={Unit cells}, answer={81}, answer_src=solved, working={fcc: $a=4r/\sqrt2=2$ \AA $=2\times10^{-8}$ cm; $M=\rho a^3N_A/Z=67.5\times8\times10^{-24}\times6\times10^{23}/4=81$ g/mol}, difficulty=4, diff_parts={K1 S2 I0 C1 F1}, diff_note={fcc a=4r/sqrt2, M=rho a^3 N_A/4}]
For any metal $^{18}$M crystallising in fcc arrangement (Atomic radius $=1/\sqrt2$ \AA) with density is 67.5 gm/cc hence atomic mass of the metal is [$N_A=6\times10^{23}$] (in gram/mole)
\end{bankq}
\begin{bankq}[id=TOAS02-P1-CHEM-S1-Q02, source={TOAS - Test 02 (13 Feb 2022)}, paper=1, qno=2, type=single, subtopic={Defects}, answer={A}, answer_src=key, difficulty=1, diff_parts={K1 S1 I0 C0 F0}, diff_note={Pure silicon is an intrinsic semiconductor}]
Which of the following is intrinsic semiconductor?
\begin{qoptions}
\opt{A}{Si}
\opt{B}{Si doped with P}
\opt{C}{GaAs}
\opt{D}{$\mathrm{Fe_{0.93}O_{1.00}}$}
\end{qoptions}
\end{bankq}
\begin{bankq}[id=TOAS02-P1-CHEM-S2-Q07, source={TOAS - Test 02 (13 Feb 2022)}, paper=1, qno=7, type=multi, subtopic={Packing}, answer={B,C,D}, answer_src=solved, working={Both are close packed with the same radius so $d_{fcc}=d_{hcp}$ (C true, A false); $a_{fcc}=2\sqrt2\,r$, $c_{hcp}=4\sqrt{2/3}\,r$ so ratio $=\sqrt3/2$ (B); $z_{fcc}/z_{hcp}=4/6=2/3$ (D)}, difficulty=6, diff_parts={K2 S3 I1 C0 F1}, diff_note={fcc and hcp densities equal; a_fcc and c_hcp in terms of r; Z ratio}]
A metal `M' crystallizes in two crystal systems fcc and hcp. If $d$-density, $z$ - number of atoms per unit cell and $a,b,c$ are edge lengths of the unit cell along the crystallographic axes then. Select the correct options :
\begin{qoptions}
\opt{A}{$\dfrac{d_{\mathrm{fcc}}}{d_{\mathrm{hcp}}}\ne1$}
\opt{B}{$\dfrac{a_{\mathrm{fcc}}}{c_{\mathrm{hcp}}}=\dfrac{\sqrt3}{2}$}
\opt{C}{$\dfrac{d_{\mathrm{fcc}}}{d_{\mathrm{hcp}}}=1$}
\opt{D}{$\dfrac{z_{\mathrm{fcc}}}{z_{\mathrm{hcp}}}=\dfrac23$}
\end{qoptions}
\end{bankq}
\begin{bankq}[id=TOAS02-P1-CHEM-S2-Q08, source={TOAS - Test 02 (13 Feb 2022)}, paper=1, qno=8, type=multi, subtopic={Unit cells}, answer={A,C,D}, answer_src=solved, working={In CaF$_2$ Ca$^{2+}$ forms ccp and F$^-$ fills all tetrahedral voids: CN of Ca$^{2+}$ = 8 (A true), CN of F$^-$ = 4 (B false), C true; equivalently Ca$^{2+}$ occupies alternate cubes of the simple cubic F$^-$ array, not all of them, so D is false}, difficulty=3, diff_parts={K1 S2 I0 C0 F1}, diff_note={Fluorite structure: coordination numbers and void occupancy}]
The correct information(s) about $\mathrm{CaF_2}$ crystal (Fluorite structure) is/are.
\begin{qoptions}
\opt{A}{C.N. of $\mathrm{Ca^{2+}}$ ion is 8}
\opt{B}{C.N. of $\mathrm{F^-}$ ion is 8}
\opt{C}{$\mathrm{F^-}$ ions occupy all the tetrahedral voids formed by CCP arrangement of $\mathrm{Ca^{2+}}$ ions}
\opt{D}{$\mathrm{Ca^{2+}}$ ions occupy all the cubic voids formed by simple cubic arrangement of $\mathrm{F^-}$ ions}
\end{qoptions}
\end{bankq}
\begin{bankq}[id=TOAS02-P1-CHEM-S3-Q01, source={TOAS - Test 02 (13 Feb 2022)}, paper=1, qno=1, type=numerical, subtopic={Unit cells}, answer={40}, answer_src=solved, working={Cell: A 6 corners = 0.75; B 2 corners + 3 mutually non-opposite face centres = 0.25+1.5; C body centre + other 3 faces = 1+1.5; mass $=0.75(20)+1.75(60)+2.5(80)=320$ amu, volume 8 \AA$^3$, density 40}, difficulty=7, diff_parts={K2 S3 I1 C1 F1}, diff_note={Atoms per cell with the given occupancy pattern, then mass over volume}]
A compound has cubical unit cell in which `A' atoms are found at six corners and `B' atoms are present at remaining corners and at face centres which are not opposite to each other and `C' atoms are present at body centres and remaining face centres. Calculate density of the compound in $\mathrm{amu/\AA^3}$.\\
Given : A = 20 amu, B = 60 amu, C = 80 amu and edge length = 2\AA.
\end{bankq}
\begin{bankq}[id=TOAS03-P1-CHEM-S1-Q04, source={TOAS - Test 03 (27 Feb 2022)}, paper=1, qno=4, type=single, subtopic={Unit cells}, answer={C}, answer_src=key, difficulty=6, diff_parts={K2 S3 I1 C1 F0}, diff_note={Count ions removed along an edge diagonal of the NaCl cell and recompute the packing fraction}]
Correct expression for determining packing fraction (P.F.) of an ideal NaCl unit cell, if all ions along an edge diagonal are absent, as shown [$\mathrm{Cl^-}$ is forming FCC, $r_+$ = Cationic radius, $r_-$ = Anionic radius]
\bfig{119.0}{TOAS03_p1_chem_s1_q04_1.png}
\begin{qoptions}
\opt{A}{$\dfrac{\frac43\pi\left(r_+^3+r_-^3\right)}{16\sqrt2\,r_-^3}$}
\opt{B}{$\dfrac{\frac43\pi\left(\frac72\cdot r_+^3+4r_-^3\right)}{16\sqrt2\,r_-^3}$}
\opt{C}{$\dfrac{\frac43\pi\left(\frac52 r_+^3+4r_-^3\right)}{16\sqrt2\,r_-^3}$}
\opt{D}{$\dfrac{\frac43\pi\left(\frac52 r_+^3+4r_-^3\right)}{16\sqrt2\,r_+^3}$}
\end{qoptions}
\end{bankq}
\begin{bankq}[id=TOAS03-P1-CHEM-S2-Q10, source={TOAS - Test 03 (27 Feb 2022)}, paper=1, qno=10, type=multi, subtopic={Defects}, answer={A,B,D}, answer_src=key, difficulty=3, diff_parts={K1 S2 I0 C0 F1}, diff_note={Four statements on defects, hexagonal cell parameters and CN of NaCl}]
Select the correct statement(s) :
\begin{qoptions}
\opt{A}{The ionic crystal of AgBr may have Schottky defect}
\opt{B}{The unit cell having crystal parameters, $a=b\neq c$, $\alpha=\beta=90^\circ$, $\gamma=120^\circ$ is hexagonal}
\opt{C}{In ionic compounds having Frenkel defect, the ratio $r^+/r^-$ is high}
\opt{D}{The co-ordination number of $\mathrm{Na^+}$ ion in NaCl is 6}
\end{qoptions}
\end{bankq}
\begin{bankq}[id=TOAS03-P1-CHEM-S3-Q04, source={TOAS - Test 03 (27 Feb 2022)}, paper=1, qno=4, type=numerical, subtopic={Unit cells}, answer={6}, answer_src=key, difficulty=3, diff_parts={K1 S2 I0 C0 F1}, diff_note={Zn-Zn nearest distance in zinc blende is a/sqrt2}]
If shortest distance between two nearest $\mathrm{Zn^{2+}}$ ion in ZnS crystal is $3\sqrt2$ \AA. Then edge length of this crystal (in \AA) will be
\end{bankq}
\begin{bankq}[id=TOAS03-P2-CHEM-S1-Q03, source={TOAS - Test 03 (27 Feb 2022)}, paper=2, qno=3, type=multi, subtopic={Packing}, answer={A,D}, answer_src=key, difficulty=6, diff_parts={K1 S4 I1 C0 F1}, diff_note={Neighbour counts and distances in cubic, antifluorite, fcc layers and ZnS}]
Select correct statement(s)
\begin{qoptions}
\opt{A}{Number of second nearest neighbours in primitive cubic unit cell is 12.}
\opt{B}{Minimum inter-ionic distance in $\mathrm{Na_2O}$ is $\sqrt3\,a/2$.}
\opt{C}{Distance between two consecutive hexagonal layer in FCC in $4r\sqrt{2/3}$.}
\opt{D}{Number of next nearest neighbours for $\mathrm{S^{2-}}$ in ZnS type structure is 12.}
\end{qoptions}
\end{bankq}
\begin{bankpassage}[id=TOAS04-P2-CHEM-S3-G3-4, source={TOAS - Test 04 (15 Mar 2022)}, questions={3-4}]
A hypothetical metal $M$ (Atomic weight = 450) crystallizes as a cubic crystal such that metal atoms touch each other along body diagonal of cube. Edge length of cubic unit cell is 1.0 nm. If this metal is melted then it's density in molten form is 1.2 gm/cc. ($N_A=6\times10^{23}$)
\end{bankpassage}
\begin{bankq}[id=TOAS04-P2-CHEM-S3-Q03, source={TOAS - Test 04 (15 Mar 2022)}, paper=2, qno=3, type=numerical, subtopic={Unit cells}, answer={1.50}, answer_src=key, difficulty=4, diff_parts={K2 S2 I0 C0 F1}, diff_note={bcc from contact along the body diagonal: density = 2M/(N_A a^3)}, group=TOAS04-P2-CHEM-S3-G3-4]
Density of solid metal in gm/cc is -
\end{bankq}
\begin{bankq}[id=TOAS04-P2-CHEM-S3-Q04, source={TOAS - Test 04 (15 Mar 2022)}, paper=2, qno=4, type=numerical, subtopic={Unit cells}, answer={45.50 to 45.70}, answer_src=key, difficulty=7, diff_parts={K2 S3 I1 C1 F1}, diff_note={Packing fraction pi sqrt3/8 gives empty space of 45.6 percent for the molten density}, group=TOAS04-P2-CHEM-S3-G3-4]
Percentage of empty space in molten state of metal will be -
\end{bankq}
\begin{bankq}[id=TOAS05-P1-CHEM-S1-Q04, source={TOAS - Test 05 (03 Apr 2022)}, paper=1, qno=4, type=multi, subtopic={Packing}, answer={B,C}, answer_src=key, difficulty=5, diff_parts={K1 S3 I1 C0 F1}, diff_note={Layers needed for primitive cubic, fcc, hcp and bcc cells}]
Choose the correct statement among the following :
\begin{qoptions}
\opt{A}{Primitive cubic unit cell involves two dimensional hexagonal closed packed array of atoms}
\opt{B}{Four hexagonal closed packed two dimensional layers of atoms are involved in one FCC unit cell}
\opt{C}{Three hexagonal close packed two dimensional layers of atoms are involved in one HCP unit cell}
\opt{D}{Four square close packed two dimensional layers of atoms are involved in one BCC unit cell}
\end{qoptions}
\end{bankq}
\begin{bankq}[id=TOAS05-P2-CHEM-S1-Q08, source={TOAS - Test 05 (03 Apr 2022)}, paper=2, qno=8, type=single, subtopic={Defects}, answer={C}, answer_src=key, difficulty=2, diff_parts={K1 S2 I0 C0 F0}, diff_note={Match magnetic property with solids (NaCl, O2, MnO, Fe3O4)}]
Match items of \underline{List--I} with those of \underline{List--II} :
{\small\setlength{\tabcolsep}{3pt}\renewcommand{\arraystretch}{1.25}\begin{center}\begin{tabular}{|L{0.6cm}|L{3.0cm}|L{0.7cm}|L{1.5cm}|}
\hline
\multicolumn{2}{|c|}{\parbox{3.4cm}{\centering List--I\\(Property)}} & \multicolumn{2}{c|}{\parbox{2.3cm}{\centering List--II\\(Example)}}\\
\hline
(a) & Diamagnetism & (i) & $\mathrm{MnO}$\\
\hline
(b) & Ferrimagnetism & (ii) & $\mathrm{O_2}$\\
\hline
(c) & Paramagnetism & (iii) & $\mathrm{NaCl}$\\
\hline
(d) & Antiferromagnetism & (iv) & $\mathrm{Fe_3O_4}$\\
\hline
\end{tabular}\end{center}}
Choose the \textbf{most appropriate} answer from the options given below :
\begin{qoptions}
\opt{A}{(a)--(ii), (b)--(i), (c)--(iii), (d)--(iv)}
\opt{B}{(a)--(iii), (b)--(i), (c)--(ii), (d)--(iv)}
\opt{C}{(a)--(iii), (b)--(iv), (c)--(ii), (d)--(i)}
\opt{D}{(a)--(iv), (b)--(ii), (c)--(i), (d)--(iii)}
\end{qoptions}
\end{bankq}
\begin{bankq}[id=TOAS05-P2-CHEM-S1-Q20, source={TOAS - Test 05 (03 Apr 2022)}, paper=2, qno=20, type=single, subtopic={Unit cells}, answer={C}, answer_src=key, difficulty=7, diff_parts={K2 S4 I2 C0 F0}, diff_note={Cut the NaCl-type cell in two identical halves: contents and lattice parameters}]
Solid $AB$ form NaCl type structure in which $A^+$ ion is present at corner \& face center of unit cell. If this unit cell is cut in two identical parts according to given plane in the figure. These two parts of unit cell called newly generated unit cells. INCORRECT option is -
\bfig{82.0}{TOAS05_p2_chem_s1_q20_1.png}
\begin{qoptions}
\opt{A}{Empirical formula of $AB$ remains unchanged in newly generated unit cell}
\opt{B}{Equal numbers of $A^+$ and $B^-$ ions are present at both edge centres and corners in newly generated unit cell}
\opt{C}{Cation $A^+$ will be present on face centres of newly generated unit cell}
\opt{D}{Newly generated unit cell will have same lattice parameters as that of tetragonal lattice system}
\end{qoptions}
\end{bankq}
\begin{bankq}[id=TOAS05-P2-CHEM-S2-Q04, source={TOAS - Test 05 (03 Apr 2022)}, paper=2, qno=4, type=numerical, subtopic={Defects}, answer={4}, answer_src=key, difficulty=7, diff_parts={K2 S3 I1 C1 F1}, diff_note={Cation vacancies from M2+ doping change the mass at constant volume}]
An ionic compound $\left(\overset{+}{X}\overset{-}{Y}\right)$ is doped with controlled amount of $(MY_2)$. Number of cation vacancies in 1 mol of $XY$ is found to be $6.02\times10^{21}$. If atomic mass of $X$, $Y$ and $M$ are 15,10 \& 20 a.m.u respectively :\\
(Assume volume remains unchanged)\\
(Given : $N_A=6.02\times10^{23}$) If the \% decreases in density is $x$, then the value of $10x$ is :
\end{bankq}
\begin{bankq}[id=TOAS06-P1-CHEM-S2-Q06, source={TOAS - Test 06 (17 Apr 2022)}, paper=1, qno=6, type=multi, subtopic={Packing}, answer={A,B}, answer_src=key, difficulty=6, diff_parts={K1 S3 I2 C0 F1}, diff_note={Cube joining the tetrahedral voids of fcc: contents are one octahedral void and 8 tetrahedral void fractions}]
In an FCC unit cell a cube is formed by joining the centers of all the tetrahedral voids to generates a new cube. Then the new cube would contain voids as :
\begin{qoptions}
\opt{A}{1 full tetrahedral void}
\opt{B}{1 full octahedral void}
\opt{C}{8 full tetrahedral void and 1 full octahedral void}
\opt{D}{1 full octahedral void only}
\end{qoptions}
\end{bankq}
\begin{bankq}[id=TOAS07-P1-CHEM-S2-Q07, source={TOAS - Test 07 (01 May 2022)}, paper=1, qno=7, type=multi, subtopic={Unit cells}, answer={A,B,D}, answer_src=key, difficulty=4, diff_parts={K1 S3 I0 C0 F1}, diff_note={Void occupancy in fluorite, zinc blende, rock salt, antifluorite}]
Which of the following statement(s) is/are correct ?
\begin{qoptions}
\opt{A}{In fluorite structure, tetrahedral voids are occupied by anions.}
\opt{B}{In zinc-blende structure, half of tetrahedral voids are occupied by cations.}
\opt{C}{In rock-salt structure, cations occupy tetrahedral voids.}
\opt{D}{The general formula of an ionic compound crystallizing in antifluorite structure is $\mathrm{A_2B}$.}
\end{qoptions}
\end{bankq}
\begin{bankq}[id=TOAS07-P2-CHEM-S1-Q03, source={TOAS - Test 07 (01 May 2022)}, paper=2, qno=3, type=multi, subtopic={Packing}, answer={D}, answer_src=key, difficulty=7, diff_parts={K2 S3 I1 C1 F1}, diff_note={fcc and bcc edge lengths, density ratio, second neighbour distance}]
Element $X$ crystallizes in a 12 co-ordination of FCC lattice. On applying high temperature it changes to 8 co-ordination BCC lattice. If $r$ is the radius of atom $X$ in both lattices, then , INCORRECT statement(s) is/are-
\begin{qoptions}
\opt{A}{Edge length of unit cell in 12 co-ordination lattice is $2\sqrt2\,r$}
\opt{B}{Edge length of unit cell in 8 co-ordination lattice is $\dfrac{4}{\sqrt3}r$}
\opt{C}{Ratio of densities before and after applying high temperature is $2(\sqrt2)^3:(\sqrt3)^3$}
\opt{D}{Distance of $2^{\mathrm{nd}}$ nearest neighbour in 8 co-ordination BCC lattice is $2\sqrt2\,r$}
\end{qoptions}
\end{bankq}
\begin{bankq}[id=TOAS08-P1-CHEM-S3-Q08, source={TOAS - Test 08 (15 May 2022)}, paper=1, qno=8, type=numerical, subtopic={Defects}, answer={58.00 to 61.00}, answer_src=key, difficulty=9, diff_parts={K2 S4 I2 C1 F1}, diff_note={Defective rock-salt cell with C in tetrahedral holes: fraction of body diagonal occupied}]
An ideal crystal $AB$ has rock salt structure in which $A^+$ is doped with atoms of `$C$' which occupy all the tetrahedral voids without distorting the lattice. If the crystal shows defect such that the body-centered atoms is missing, then the percentage of body diagonal covered by ions/atoms in the defective crystal is (Given $\sqrt{1.5}=29/24$)
\end{bankq}
\begin{bankq}[id=TOAS08-P2-CHEM-S2-Q08, source={TOAS - Test 08 (15 May 2022)}, paper=2, qno=8, type=single, subtopic={Unit cells}, answer={A}, answer_src=key, difficulty=4, diff_parts={K1 S3 I1 C0 F0}, diff_note={Match structures with sublattice occupations}]
Match the List.
\begin{center}\footnotesize\renewcommand{\arraystretch}{1.1}\setlength{\extrarowheight}{3.5pt}\setlength{\tabcolsep}{2pt}
\begin{tabular}{|c|l|c|l|}
\hline
\multicolumn{2}{|c|}{List-I} & \multicolumn{2}{c|}{List-II}\\[3pt]
\hline
(P) & NaCl structure & (1) & \parbox[t]{3.1cm}{Cation forming fcc and anion at octahedral void}\\[3pt]
(Q) & Zinc Blende structure & (2) & \parbox[t]{3.1cm}{Anion forming fcc and at all octahedral void while cation goes to all tetrahedral void}\\[3pt]
(R) & Ceasium Halide structure & (3) & \parbox[t]{3.1cm}{Anion forming fcc and cation at alternate tetrahedral void}\\[3pt]
(S) & Flourite structure & (4) & \parbox[t]{3.1cm}{Cation forming fcc and anion at alternate tetrahedral void.}\\[3pt]
 & & (5) & \parbox[t]{3.1cm}{Anion forming simple cube \& cation at cubical void.}\\[3pt]
 & & (6) & \parbox[t]{3.1cm}{Cation forming fcc \& anion at all tetrahedral void.}\\[3pt]
\hline
\end{tabular}
\end{center}
Which of the following option is correct match?
\begin{qoptions}
\opt{A}{$P\to1;Q\to3,4;R\to2,5;S\to6$}
\opt{B}{$P\to1;Q\to3,4;R\to2;S\to6$}
\opt{C}{$P\to3,4;Q\to1;R\to2;S\to6$}
\opt{D}{$P\to3;Q\to4;R\to2,5;S\to1$}
\end{qoptions}
\end{bankq}
\begin{bankq}[id=TOAS08-P2-CHEM-S3-Q04, source={TOAS - Test 08 (15 May 2022)}, paper=2, qno=4, type=numerical, subtopic={Packing}, answer={1.00}, answer_src=key, difficulty=9, diff_parts={K2 S4 I2 C1 F1}, diff_note={Boiling-point elevation gives moles of ions; cell volume and radius ratio}]
A solid cube with edge length of 2.532 cm of an ionic compound which has NaCl type lattice is added to 1kg of water. The boiling point of this solution is found to be $100.52^\circ$C ( Assume 100\% ionisation of ionic compound). If radius of anion of ionic solid is 2\,\AA{}, then calculate radius of cation in \AA{}.\newline
Given : $k_b$ of water $=0.52$ K kg mol$^{-1}$,\\[2pt]
$N_A=6\times10^{23}$, $\sqrt[3]{216}=6$, $(2.532)^3\times\dfrac43=21.6$
\end{bankq}
\begin{bankq}[id=TOAS08-P2-CHEM-S3-Q05, source={TOAS - Test 08 (15 May 2022)}, paper=2, qno=5, type=numerical, subtopic={Unit cells}, answer={3.00}, answer_src=key, difficulty=6, diff_parts={K1 S3 I1 C1 F1}, diff_note={Distances in the fluorite lattice and the expression}]
In $\mathrm{CaF_2}$-type lattice, if\newline
$x=$ nearest distance between $\mathrm{Ca^{2+}}$ \& $\mathrm{F^-}$ ions\newline
$y=$ nearest distance between two $\mathrm{F^-}$ ions\newline
$z=$ nearest distance between two $\mathrm{Ca^{2+}}$ ions.\newline
then find value of the expression.
\[\left[\sqrt6\left(\dfrac{xz}{y^2}\right)\right]\]
\end{bankq}
\begin{bankq}[id=TOAS09-P1-CHEM-S1-Q03, source={TOAS - Test 09 (30 May 2022)}, paper=1, qno=3, type=single, subtopic={Defects}, answer={C}, answer_src=key, difficulty=1, diff_parts={K1 S1 I0 C0 F0}, diff_note={Defects: standard approach for this subtopic; 1 concept(s), 1 step(s)}]
Select the incorrect statement :
\begin{qoptions}
\opt{A}{MgFe$_2$O$_4$ can behave like a ferrimagnetic solid}
\opt{B}{The unit cell having crystal parameters, $a=b\neq c$, $\alpha=\beta=90^\circ$, $\gamma=120^\circ$ is hexagonal}
\opt{C}{In ionic compounds having Frenkel defect, the ratio $r^+/r^-$ is high}
\opt{D}{The co-ordination number of Na$^+$ ion in NaCl is 6}
\end{qoptions}
\end{bankq}
\begin{bankq}[id=TOAS09-P1-CHEM-S1-Q04, source={TOAS - Test 09 (30 May 2022)}, paper=1, qno=4, type=single, subtopic={Defects}, answer={B}, answer_src=key, difficulty=1, diff_parts={K1 S1 I0 C0 F0}, diff_note={Defects: standard approach for this subtopic; 1 concept(s), 1 step(s)}]
If FeO crystallizes rock salt structure in which O$^{2-}$ ions occupy FCC lattice points. Due to some crystal defect, lattice becomes Fe$_{0.75}$O. Find number of effective Fe$^{3+}$ ions per unit cell -
\begin{qoptions}
\opt{A}{1}
\opt{B}{2}
\opt{C}{3}
\opt{D}{4}
\end{qoptions}
\end{bankq}
\begin{bankq}[id=TOAS09-P2-CHEM-S1-Q20, source={TOAS - Test 09 (30 May 2022)}, paper=2, qno=20, type=single, subtopic={Unit cells}, answer={B}, answer_src=key, difficulty=4, diff_parts={K2 S2 I1 C0 F0}, diff_note={Unit cells: standard approach for this subtopic; 2 concept(s), 2 step(s), needs a non-obvious insight}]
In a CCP lattice of X \& Y, X atoms are present at the corners while Y atoms are at face centers. Then the formula of the compound would be if one of the X-atoms from a corner in each unit cell is replaced by Z-atom :
\begin{qoptions}
\opt{A}{X$_7$ Y$_{24}$ Z$_2$}
\opt{B}{X$_7$ Y$_{24}$ Z}
\opt{C}{X$_{24}$ Y$_7$ Z}
\opt{D}{X Y$_{24}$ Z}
\end{qoptions}
\end{bankq}
\begin{bankq}[id=TOAS09-P2-CHEM-S2-Q10, source={TOAS - Test 09 (30 May 2022)}, paper=2, qno=10, type=numerical, subtopic={Unit cells}, answer={2}, answer_src=key, difficulty=7, diff_parts={K2 S3 I1 C1 F1}, diff_note={Unit cells: standard approach for this subtopic; 2 concept(s), 3 step(s), needs a non-obvious insight, heavier computation}]
A metal `A' exists as BCC structure having uncovered length between the adjacent atoms along the edge in an unit cell equal to 1.2 \AA{}. If molar mass of the metal is 307.2 gm/mole then calculate density of the crystal in gm/ml. [Use : $N_A=6\times10^{23}$ and $\sqrt3=1.7$]
\end{bankq}
\end{document}
